\chapter{LAMPIRAN A}

Lampiran ini disiapkan untuk artefak yang akan dipublikasikan setelah protokol ditinjau, apparatus divalidasi, dan pengumpulan data formal dilakukan. Pada draf ini, skrip, konfigurasi aktual, dan dataset tervalidasi belum dilampirkan. Ketiadaan artefak tersebut tidak menunjukkan bahwa pengujian telah berhasil atau gagal.

\section{Artefak Eksperimen yang Direncanakan}

Lampiran versi akhir dapat memuat sumber pemeriksa dan korektor imperatif berbasis \texttt{kubectl}, konfigurasi aplikasi ArgoCD, manifest dan inventaris workload formal, skema kolektor, serta prosedur pemulihan nilai dasar. Setiap artefak perlu diberi versi dan dikaitkan dengan konfigurasi aktual eksperimen. Materi pilot S1/S2 harus dipisahkan dari artefak dan data pengujian formal.

\section{Skema Data dan Perlindungan Informasi}

Skema log formal perlu menyimpan pengenal run, lingkungan, skenario pada berkas metadata terpisah dari masukan pengendali, status penerimaan injeksi, cap waktu monotonik dalam detik dan UTC untuk $t_0$, $t_a$, $t_d$, $t_e$, dan $t_r$, interval polling aktual, hasil kesetaraan dan kesehatan, status penyensoran, kesalahan, serta alasan run tidak valid. $t_a$ menyimpan waktu penerimaan respons API injeksi. Metadata pengenal skenario tidak boleh diberikan kepada mekanisme pendeteksi atau korektor.

Snapshot dan bukti harus disanitasi sebelum disimpan atau dibagikan. Nilai Secret tidak dimasukkan ke listing, tangkapan layar, log, atau dataset publik. Pencocokan Secret dilakukan di memori; artefak yang dibagikan hanya memuat identitas bidang dan hasil kesetaraan, tanpa nilai rahasia atau hash yang dapat dipakai untuk menebak isinya.

\section{Status Lampiran}

Seluruh item pada lampiran ini masih berupa kebutuhan artefak, bukan skrip atau konfigurasi yang telah diverifikasi. Isi lampiran hanya dapat dilengkapi berdasarkan bukti yang benar-benar dihasilkan oleh eksperimen dan telah ditinjau untuk mencegah kebocoran informasi sensitif.
